\section{Discrete reference implementation (appendix)}
\label{app:discrete}

This appendix records a \emph{reference} discrete implementation intended to approximate the continuum model.
It is not part of the conceptual claims in the main text.

\subsection{Purpose and philosophy}

The discrete model is used to:
\begin{itemize}
\item test isotropy numerically (including diagonal couplings as an approximation to rotational invariance),
\item extract dispersion relations for small-amplitude waves,
\item compute geometric-phase diagnostics from simulated polarized wave packets,
\item search for and stress-test localized modes.
\end{itemize}

\subsection{Lattice kinematics (summary)}

A cubic lattice of nodes approximates $\Omega\subset\R^3$. Each node carries a 4D position $\vecX_{ijk}(t)\in\R^4$.
Springs connect near neighbors; including diagonals improves isotropy.
A uniform prestretch parameter $\alpha$ sets the spring rest lengths relative to the lattice spacing,
representing homogeneous pre-tension.

\subsection{Time integration and invariants}

A symplectic or near-symplectic second-order integrator (e.g.\ Velocity Verlet) is used.
Conservation checks include total energy drift and momentum balance, with attention to boundary conditions.

\subsection{Geometric-phase extraction in the discrete setting}

Given a narrowband simulated wave, a local polarization state $\bm{\epsilon}$ can be estimated from
phase-aligned projections onto an internal basis (e.g.\ transverse directions).
Then $A_\mu$ and $F_{\mu\nu}$ are computed by discrete derivatives, and loop holonomies are evaluated by discrete line integrals.

\subsection{Recommended reporting}

To reduce ambiguity, report:
grid size, timestep, boundary conditions, constitutive/spring law, prestretch $\alpha$,
measured dispersion curves, and numerical stability metrics alongside any holonomy results.
